\documentclass[12pt,a4paper]{article}
\usepackage{ctex}
\title{基于ViT的CIFAR-10图像分类实验报告}
\author{}

\usepackage{braket}
\usepackage{hyperref}
\hypersetup{colorlinks=true, linkcolor=blue, anchorcolor=green, citecolor=green}
\usepackage{bm}
\usepackage{amssymb, amsmath, amsthm}
\usepackage{tikz}
\usepackage{graphicx}
\usepackage{float}
\usepackage{booktabs}
\usepackage{multirow}
\usepackage{subcaption}
\usetikzlibrary{arrows.meta}

\newtheorem{theorem}{定理}[section]
\newtheorem{definition}{定义}[section]

\begin{document}

\begin{titlepage}
    \centering
    \vspace*{1cm}
    
\includegraphics[width=0.6\textwidth]{ucas_logo.pdf}
    \vspace{1cm}
    
    {\Huge\bfseries 深度学习作业2\par}
    \vspace{1cm}
    {\LARGE  基于ViT的CIFAR-10图像分类实验\par}
    \vspace{3cm}

    \begin{tabular}{rl}
        {\kaishu\Large 学号：} & {\kaishu\Large \underline{\makebox[6cm]{202528014629014}}} \\[0.8cm]
        {\kaishu\Large 姓名：} & {\kaishu\Large \underline{\makebox[6cm]{刘浩}}} \\[0.8cm]
        {\kaishu\Large 院所：} & {\kaishu\Large \underline{\makebox[6cm]{人工智能学院}}} \\[0.8cm]
    \end{tabular}

    \vfill

    {\large \today\par}
\end{titlepage}

% \begin{titlepage}
%     \centering
%     \vspace*{2cm}

%     {\Huge\bfseries 深度学习作业2\par}
%     \vspace{1cm}
%     {\huge 基于ViT的CIFAR-10图像分类实验\par}
%     \vspace{3cm}

%     \begin{tabular}{rl}
%         {\kaishu\Large 学号：} & {\kaishu\Large \underline{\makebox[6cm]{202528014629014}}} \\[0.8cm]
%         {\kaishu\Large 姓名：} & {\kaishu\Large \underline{\makebox[6cm]{刘浩}}} \\[0.8cm]
%         {\kaishu\Large 院所：} & {\kaishu\Large \underline{\makebox[6cm]{人工智能学院}}} \\[0.8cm]
%     \end{tabular}

%     \vfill

%     {\large \today\par}
% \end{titlepage}

\newpage
\tableofcontents
\newpage

% ==============================================================
\section{实验目的}
% ==============================================================

\begin{enumerate}
    \item 理解Vision Transformer (ViT) 的核心思想，掌握自注意力机制在图像分类中的应用
    \item 从零实现ViT模型，包括Patch Embedding、位置编码、Transformer Encoder等关键组件
    \item 在CIFAR-10数据集上训练ViT模型，达到80\%以上的测试准确率
    \item 深入理解模型的训练行为与分类特性
\end{enumerate}

% ==============================================================
\section{实验原理}
% ==============================================================

\subsection{Vision Transformer 架构概述}

Vision Transformer (ViT) 由Dosovitskiy等人于2020年提出\cite{dosovitskiy2020image}，其核心思想是将图像视为一组视觉"词元"（token），直接应用NLP领域的Transformer架构进行图像分类。与CNN依赖局部卷积不同，ViT通过全局自注意力机制捕获图像各区域之间的长程依赖关系。

ViT的整体流程为：
\begin{equation}
\begin{aligned}
\text{Image} &\xrightarrow{\text{Patch Embed}} \text{Token序列} \\
             &\xrightarrow{+\text{Pos Embed}} \text{Transformer Encoder} \\
             &\xrightarrow{\text{CLS Token}} \text{MLP Head} \\
             &\to \text{类别预测}
\end{aligned}
\end{equation}

\subsection{Patch Embedding}

给定输入图像 $\bm{x} \in \mathbb{R}^{H \times W \times C}$，ViT首先将其切分为大小为 $P \times P$ 的非重叠图像块（patch）。对于 $32 \times 32$ 的CIFAR-10图像、$P=4$ 的patch大小，共产生 $N = (32/4)^2 = 64$ 个patch。

每个patch被展平为向量 $\bm{x}_p^i \in \mathbb{R}^{P^2 \cdot C}$（维度为 $4 \times 4 \times 3 = 48$），再通过线性投影映射到 $D$ 维嵌入空间：
\begin{equation}
\bm{z}_0^i = \bm{E} \bm{x}_p^i + \bm{e}_{pos}^i, \quad \bm{E} \in \mathbb{R}^{D \times (P^2 C)}
\end{equation}

其中 $\bm{e}_{pos}^i$ 为可学习的位置编码，用于保留空间位置信息。此外，在序列前端拼接一个可学习的分类token $\bm{z}_0^{[\text{CLS}]}$，使总序列长度为 $N+1=65$。

\subsection{多头自注意力机制}

Transformer Encoder的核心是多头自注意力机制（Multi-Head Self-Attention, MHSA）。对输入 $\bm{Z} \in \mathbb{R}^{(N+1) \times D}$，首先通过线性投影得到查询、键、值矩阵：
\begin{equation}
\bm{Q} = \bm{Z}\bm{W}_Q, \quad \bm{K} = \bm{Z}\bm{W}_K, \quad \bm{V} = \bm{Z}\bm{W}_V
\end{equation}

其中 $\bm{W}_Q, \bm{W}_K, \bm{W}_V \in \mathbb{R}^{D \times d_k}$。注意力计算为：
\begin{equation}
\text{Attention}(\bm{Q}, \bm{K}, \bm{V}) = \text{softmax}\left(\frac{\bm{Q}\bm{K}^\top}{\sqrt{d_k}}\right)\bm{V}
\end{equation}

多头注意力将嵌入维度拆分为 $h$ 个头，每个头独立计算注意力后拼接：
\begin{equation}
\text{MHSA}(\bm{Z}) = \text{Concat}(\text{head}_1, \dots, \text{head}_h)\bm{W}_O
\end{equation}

本实验使用 $h=8$ 个注意力头，每个头的维度 $d_k = 64$。

\subsection{Transformer Encoder Block}

每个Transformer Block由MHSA和前馈网络（FFN）组成，均采用Pre-Norm残差连接：
\begin{align}
\bm{z}'_l &= \text{MHSA}(\text{LN}(\bm{z}_{l-1})) + \bm{z}_{l-1} \\
\bm{z}_l &= \text{FFN}(\text{LN}(\bm{z}'_l)) + \bm{z}'_l
\end{align}

FFN为两层全连接网络，使用GELU激活函数：
\begin{equation}
\text{FFN}(\bm{x}) = \bm{W}_2 \cdot \text{GELU}(\bm{W}_1 \bm{x} + \bm{b}_1) + \bm{b}_2
\end{equation}

\subsection{分类头}

经过 $L$ 层Transformer Encoder后，取CLS token的输出 $\bm{z}_L^{[\text{CLS}]}$，通过线性层映射到类别数：
\begin{equation}
\hat{y} = \text{Linear}(\text{LN}(\bm{z}_L^{[\text{CLS}]})) \in \mathbb{R}^{10}
\end{equation}

% ==============================================================
\section{实验环境与数据集}
% ==============================================================

\subsection{实验环境}

\begin{table}[H]
\centering
\begin{tabular}{cc}
\toprule
工具/框架 & 说明 \\
\midrule
Python 3.10 & 编程语言 \\
PyTorch & 深度学习框架 \\
CUDA & GPU加速计算 \\
einops & 张量重排工具 \\
Matplotlib / Seaborn & 数据可视化 \\
scikit-learn & 评估指标计算 \\
\bottomrule
\end{tabular}
\caption{实验环境配置}
\label{tab:environment}
\end{table}

\subsection{数据集}

CIFAR-10数据集包含60000张$32 \times 32$像素的彩色图像，分为10个类别：飞机（Airplane）、汽车（Automobile）、鸟（Bird）、猫（Cat）、鹿（Deer）、狗（Dog）、蛙（Frog）、马（Horse）、船（Ship）和卡车（Truck）。其中50000张用于训练，10000张用于测试，每个类别各有1000张测试图像。

相比实验一中使用的MNIST数据集，CIFAR-10具有以下挑战：
\begin{itemize}
    \item \textbf{彩色图像}：三通道RGB图像，信息量更大
    \item \textbf{类内多样性高}：同一类别的物体在姿态、光照、遮挡等方面差异显著
    \item \textbf{类间相似性}：如猫与狗、卡车与汽车在视觉特征上高度相似
\end{itemize}

数据增强策略包括随机裁剪和随机水平翻转，并使用ImageNet均值和标准差进行归一化。

% ==============================================================
\section{模型设计与实现}
% ==============================================================

\subsection{模型架构}

本实验从零实现了一个适配CIFAR-10的小型ViT模型。由于CIFAR-10图像尺寸仅为$32 \times 32$（远小于ViT原论文中的$224 \times 224$），我们使用较小的patch大小和模型维度以避免参数冗余。

\begin{table}[H]
\centering
\begin{tabular}{ccc}
\toprule
超参数 & 值 & 说明 \\
\midrule
Image Size & 32 & CIFAR-10图像尺寸 \\
Patch Size & 4 & 每个patch为$4\times 4$像素 \\
Embedding Dim ($D$) & 256 & token嵌入维度 \\
Depth ($L$) & 6 & Transformer层数 \\
Heads ($h$) & 8 & 注意力头数 \\
Head Dim ($d_k$) & 64 & 每头维度 \\
MLP Dim & 512 & FFN隐藏层维度 \\
Dropout & 0.1 & Dropout率 \\
\bottomrule
\end{tabular}
\caption{ViT模型超参数}
\label{tab:model_config}
\end{table}

模型的特征维度变化如下：
\begin{equation*}
\underbrace{3 \times 32 \times 32}_{\text{输入图像}} \to \underbrace{64 \times 48}_{\text{展平patch}} \to \underbrace{65 \times 256}_{\text{+CLS token}} \xrightarrow{\times 6 \text{ layers}} \underbrace{65 \times 256}_{\text{Transformer输出}} \to \underbrace{256}_{\text{CLS}} \to \underbrace{10}_{\text{类别}}
\end{equation*}

模型总参数量为\textbf{4.76M}。

\subsection{训练配置}

\begin{table}[H]
\centering
\begin{tabular}{ccc}
\toprule
参数 & 配置 & 说明 \\
\midrule
优化器 & AdamW & 带权重衰减的自适应优化器 \\
学习率 & $3 \times 10^{-4}$ & 初始学习率 \\
权重衰减 & 0.05 & 正则化强度 \\
学习率调度 & CosineAnnealing & 余弦退火学习率衰减 \\
损失函数 & CrossEntropyLoss & 多分类交叉熵损失 \\
批大小 & 128 & 每批处理样本数 \\
总Epoch & 80 & 训练轮数 \\
设备 & NVIDIA-A100 & 使用GPU加速 \\
\bottomrule
\end{tabular}
\caption{训练配置参数}
\label{tab:training_config}
\end{table}

训练采用标准流程：前向传播计算交叉熵损失，反向传播更新参数，每个epoch后在测试集评估性能。使用余弦退火学习率调度器（图\ref{fig:lr}），使学习率从$3 \times 10^{-4}$平滑衰减至接近0，有助于模型在训练后期精细调优。

% ==============================================================
\section{实验结果与分析}
% ==============================================================

\subsection{总体性能}

模型在第70个epoch达到最优测试准确率\textbf{83.29\%}，成功超过80\%的目标要求。总体指标如表\ref{tab:results}所示，宏平均Precision、Recall与F1-Score均在83\%左右，说明模型未严重偏向特定类别。

\begin{table}[H]
\centering
\begin{tabular}{cc}
\toprule
指标 & 结果 \\
\midrule
最优测试准确率 & \textbf{83.29\%} \\
最优准确率对应epoch & 第70个 \\
最终训练准确率 & 93.36\% \\
模型参数量 & 4.76M \\
Macro F1-Score & 83.13\% \\
\bottomrule
\end{tabular}
\caption{实验结果统计}
\label{tab:results}
\end{table}

\subsection{训练收敛过程}

为了理解模型如何达到上述性能，我们首先考察训练过程中的损失与准确率变化（图\ref{fig:loss}、\ref{fig:accuracy}）。

\begin{figure}[H]
\centering
\includegraphics[width=0.85\textwidth]{fig1_loss.png}
\caption{训练与验证损失变化曲线（含EMA平滑）}
\label{fig:loss}
\end{figure}

\begin{figure}[H]
\centering
\includegraphics[width=0.85\textwidth]{fig2_accuracy.png}
\caption{训练与验证准确率变化曲线}
\label{fig:accuracy}
\end{figure}

从两张曲线中可以辨认出三个阶段：\textbf{（1）快速学习期}（epoch 1--30）：训练损失从1.597急剧下降至0.48，测试准确率从52.6\%攀升至约79\%，模型在此阶段学习到了区分大类的基本视觉特征；\textbf{（2）精细调优期}（epoch 30--55）：损失下降速度放缓，测试准确率以每epoch约0.1\%的速度缓慢提升，模型在第33个epoch首次达到80\%的目标线；\textbf{（3）饱和期}（epoch 55--80）：测试准确率在83\%附近波动，验证损失停滞在0.60左右，而训练损失仍在下降——训练集与测试集的表现开始分离。

这种分离在泛化间隙图（图\ref{fig:gap}）中体现得更为清晰：前20个epoch间隙不足5\%，40个epoch后突破10\%。这是ViT在小数据集上的典型表现——Transformer缺乏CNN的归纳偏置（局部性、平移等变性），需要更大规模的数据才能充分泛化。不过，得益于Dropout（0.1）和权重衰减（0.05）的正则化，测试准确率并未出现下降，过拟合被控制在了可接受的范围内。

\begin{figure}[H]
\centering
\includegraphics[width=0.85\textwidth]{fig3_gap.png}
\caption{泛化间隙随训练进程的变化}
\label{fig:gap}
\end{figure}

图\ref{fig:dynamics}进一步从"损失改善速率"和"Train-Val损失轨迹"两个角度验证了上述三阶段划分。左图中，训练初期的改善速率高达0.04/epoch，在40个epoch后验证损失的改善速率已降为负值（即验证损失开始回升）；右图中，散点从$y=x$线附近逐渐偏离到上方，直观地反映了训练损失与验证损失的"剪刀差"。

\begin{figure}[H]
\centering
\includegraphics[width=0.95\textwidth]{fig11_dynamics.png}
\caption{训练动态：(左) 损失改善速率 (右) Train-Val损失轨迹散点图}
\label{fig:dynamics}
\end{figure}

\subsection{分类性能与错误模式}

模型整体达到了83.29\%的准确率，但各类别的表现差异显著。图\ref{fig:per_class}将10个类别按准确率升序排列：Automobile（94.0\%）和Frog（91.3\%）位列前两名，而Cat（64.3\%）和Dog（72.8\%）明显落后，Cat甚至是唯一未达到80\%目标线的类别。这一分布揭示了一个清晰的规律：\textbf{人造物体}（Automobile、Ship、Truck）具有规则的几何轮廓和鲜明的纹理，ViT的全局注意力能够轻松捕获；而\textbf{动物类别}（Cat、Dog、Bird）形态多变、类间相似度高，对细粒度判别能力要求更高。

\begin{figure}[H]
\centering
\includegraphics[width=0.85\textwidth]{fig6_per_class.png}
\caption{各类别测试准确率}
\label{fig:per_class}
\end{figure}

那么，这些低准确率类别的错误究竟发生在哪里？混淆矩阵（图\ref{fig:confusion}）给出了答案。对角线上大部分类别的正确率在80\%以上，但Cat行中有13.8\%被误判为Dog、5.4\%被误判为Bird、4.9\%被误判为Frog——Cat的错误几乎"分散"到了所有其他动物类别。反过来，Dog行中也有10.7\%被误判为Cat。Cat与Dog之间的双向混淆合计245次，占全部1671个错误的14.7\%，是影响整体准确率的最大瓶颈。

\begin{figure}[H]
\centering
\includegraphics[width=0.85\textwidth]{fig5_confusion.png}
\caption{归一化混淆矩阵}
\label{fig:confusion}
\end{figure}

为了更系统地展示类别间的语义关系，图\ref{fig:semantic}对混淆矩阵进行了对称化处理，得到类间"混淆亲和度"热力图。从中可以清晰地辨识出三个语义混淆簇：\textbf{Cat--Dog}（四足哺乳动物，毛皮纹理相似）、\textbf{Automobile--Truck}（道路交通工具，共享车轮车身特征）、以及\textbf{Bird--Deer--Frog}（自然背景中的中小型动物）。这些混淆簇与人类直觉高度一致，表明模型的错误并非随机，而是集中在视觉语义接近的类别对上。

\begin{figure}[H]
\centering
\includegraphics[width=0.85\textwidth]{fig12_semantic.png}
\caption{类间语义混淆亲和热力图}
\label{fig:semantic}
\end{figure}

图\ref{fig:samples}展示了部分正确和错误分类的样本。正确样本通常主体清晰、特征典型；而错误样本往往存在视觉歧义——例如低分辨率下某些猫的图像确实与狗难以区分，这进一步佐证了上述语义混淆的分析。

\begin{figure}[H]
\centering
\includegraphics[width=0.95\textwidth]{fig9_samples.png}
\caption{正确与错误分类的典型样本}
\label{fig:samples}
\end{figure}

最后，图\ref{fig:radar}以雷达图综合展示各类别的准确率与F1分数。两条曲线形态高度一致，说明各类别的Precision和Recall较为均衡，不存在某个类别"精确但遗漏严重"的极端情况。Cat的明显凹陷与Automobile/Frog的凸起，概括了本模型在CIFAR-10上的优势与短板。

\begin{figure}[H]
\centering
\includegraphics[width=0.6\textwidth]{fig10_radar.png}
\caption{各类别准确率与F1分数雷达图}
\label{fig:radar}
\end{figure}

% ==============================================================
\section{总结与思考}
% ==============================================================

\subsection{实验总结}

本实验从零实现了Vision Transformer模型，并在CIFAR-10数据集上达到了\textbf{83.29\%}的最优测试准确率，成功超过80\%的目标要求。主要发现如下：

\begin{enumerate}
    \item \textbf{ViT在小数据集上的局限性}：与CNN相比，ViT缺乏归纳偏置（局部性、平移等变性），在仅50000张训练图像的CIFAR-10上出现了约10\%的泛化间隙。这符合ViT原论文的结论——Transformer需要大规模数据集（如ImageNet-21k或JFT-300M）才能充分发挥优势。

    \item \textbf{类别间的语义混淆}：猫与狗之间的互相误判是影响整体准确率的最大因素（245次/1671次总错误），这揭示了$32 \times 32$低分辨率下细粒度分类的固有困难。

    \item \textbf{人造物体vs自然物体}：模型对几何结构规则的人造物体（Automobile: 94.0\%, Ship: 87.7\%）的分类效果显著优于形态多变的动物类别（Cat: 64.3\%, Dog: 72.8\%），说明ViT的全局注意力更擅长捕获规则化的视觉模式。

    \item \textbf{训练策略的有效性}：CosineAnnealing学习率调度、AdamW优化器和适度的正则化（Dropout + Weight Decay）组合在ViT训练中表现稳定，损失曲线平滑无振荡。
\end{enumerate}

\subsection{可能的改进方向}

\begin{itemize}
    \item \textbf{数据增强增强}：引入更强的数据增强（如MixUp、CutMix、RandAugment）可以有效缓解小数据集上的过拟合问题
    \item \textbf{混合架构}：使用CNN提取局部特征后再输入Transformer（如CvT、CoAtNet），结合两者的优势
    \item \textbf{知识蒸馏}：使用预训练的大型模型蒸馏知识到小型ViT，提升小模型的泛化能力（如DeiT的方法）
    \item \textbf{位置编码改进}：采用2D相对位置编码替代可学习的绝对位置编码，更好地利用图像的空间结构
\end{itemize}

\begin{thebibliography}{9}
\bibitem{dosovitskiy2020image} Dosovitskiy, A., et al. An Image is Worth 16x16 Words: Transformers for Image Recognition at Scale. \textit{ICLR}, 2021.
\end{thebibliography}

\end{document}
